\section{Markov Chain Monte Carlo - MCMC}

\subsection{Por quê Precisamos de MCMC?}

\subsection{Algoritmos de MCMC}
% Analogia da Montanha TikZ para regra de aceitação/rejeição

\subsubsection{Metropolis}
\begin{frame}[fragile]{Algoritmo de Metropolis\footnote{\citet{metropolisEquationStateCalculations1953}}}
    \footnotesize

    \SetAlCapFnt{\footnotesize}
    \SetAlCapNameFnt{\footnotesize}
    \begin{algorithm}[H]
    \DontPrintSemicolon
    \SetAlgoNoEnd
    \SetAlgoLined
    Defina um ponto inicial $\boldsymbol{\theta}^{(0)} \in \mathbb{R}^p$ do qual $P\left(\boldsymbol{\theta}^{(0)} \mid y \right) > 0$\;
     \Para{$t = 1, 2, \dots$}{
      Amostra uma proposta $\boldsymbol{\theta}^{(*)}$ de uma distribuição de propostas no tempo $t$, $J_t \left(\boldsymbol{\theta}^{(*)} \mid \boldsymbol{\theta}^{(t-1)} \right)$\;
      Como regra de aceitação/rejeição calcule a proporção das probabilidades:
      $r = \frac{P\left(\boldsymbol{\theta}^{(*)}  \mid y \right)}{P\left(\boldsymbol{\theta}^{(t-1)} \mid y \right)}$\;
      Designe:
      $
        \boldsymbol{\theta}^{(t)} =
          \begin{cases}
          \boldsymbol{\theta}^{(*)} & \text{com probabilidade $\min(r,1)$}\\
          \boldsymbol{\theta}^{(t-1)} & \text{caso contrário}
        \end{cases}
      $\;
     }
     \caption{Metropolis}
    \end{algorithm}
\end{frame}

\subsubsection{Metropolis-Hastings}
\begin{frame}[fragile]{Algoritmo de Metropolis - Hastings\footnote{\citet{hastingsMonteCarloSampling1970}}}
    \footnotesize

    \SetAlCapFnt{\footnotesize}
    \SetAlCapNameFnt{\footnotesize}
    \begin{algorithm}[H]
    \DontPrintSemicolon
    \SetAlgoNoEnd
    \SetAlgoLined
    Defina um ponto inicial $\boldsymbol{\theta}^{(0)} \in \mathbb{R}^p$ do qual $P\left(\boldsymbol{\theta}^{(0)} \mid y \right) > 0$\;
     \Para{$t = 1, 2, \dots$}{
      Amostra uma proposta $\boldsymbol{\theta}^{(*)}$ de uma distribuição de propostas no tempo $t$, $J_t \left(\boldsymbol{\theta}^{(*)} \mid \boldsymbol{\theta}^{(t-1)} \right)$\;
      Como regra de aceitação/rejeição calcule a proporção das probabilidades:
      $r = \frac{\frac{P \left(\boldsymbol{\theta}^{(*)} \mid y \right)}{J_t \left(\boldsymbol{\theta}^{(*)} \mid \boldsymbol{\theta}^{(t-1)} \right)}}{\frac{P \left(\boldsymbol{\theta}^{(t-1)} \mid y \right)}{J_t \left(\boldsymbol{\theta}^{(t-1)} \mid \boldsymbol{\theta}^{(*)} \right)}}$\;
      Designe:
      $
        \boldsymbol{\theta}^{(t)} =
          \begin{cases}
          \boldsymbol{\theta}^{(*)} & \text{com probabilidade $\min(r,1)$}\\
          \boldsymbol{\theta}^{(t-1)} & \text{caso contrário}
        \end{cases}
      $\;
     }
     \caption{Metropolis-Hastings}
    \end{algorithm}
\end{frame}

\subsubsubsection{As Limitações de Metropolis e Metropolis-Hastings}

\subsubsection{Gibbs}
\begin{frame}[fragile]{Algoritmo de Gibbs\footnote{\citet{gemanStochasticRelaxationGibbs1984}}}
    \footnotesize

    \SetAlCapFnt{\footnotesize}
    \SetAlCapNameFnt{\footnotesize}
    \begin{algorithm}[H]
    \DontPrintSemicolon
    \SetAlgoNoEnd
    \SetAlgoLined
    Defina um ponto inicial $\boldsymbol{\theta}^{(0)} \in \mathbb{R}^p$ do qual $P\left(\boldsymbol{\theta}^{(0)} \mid y \right) > 0$\;
     \Para{$t = 1, 2, \dots$}{
      Designe:
      $ \boldsymbol{\theta}^{(t)} =
        \begin{cases}
        \theta^{(t)}_1 &\sim P \left(\theta_1 \mid \theta^{(0)}_2, \dots, \theta^{(0)}_p \right) \\
        \theta^{(t)}_2 &\sim P \left(\theta_2 \mid \theta^{(t-1)}_1, \dots, \theta^{(0)}_p \right) \\
        &\vdots \\
        \theta^{(t)}_p &\sim P \left(\theta_p \mid \theta^{(t-1)}_1, \dots, \theta^{(t-1)}_{p-1} \right)
     \end{cases}
      $\;
     }
     \caption{Gibbs}
    \end{algorithm}
\end{frame}

\subsubsubsection{As Limitações de Gibbs}

\subsubsection{Hamiltonian Monte Carlo (HMC)}
\begin{frame}[fragile]{Algoritmo de HMC\footnote{\citet{neal2011mcmc}}}
    \footnotesize
    
    \SetAlCapFnt{\footnotesize}
    \SetAlCapNameFnt{\footnotesize}
    \begin{algorithm}[H]
    \DontPrintSemicolon
    \SetAlgoNoEnd
    \SetAlgoLined
    Defina um ponto inicial $\boldsymbol{\theta}^{(0)} \in \mathbb{R}^p$ do qual $P\left(\boldsymbol{\theta}^{(0)} \mid y \right) > 0$\;
    Amostre $\boldsymbol{\phi}$ de uma $\text{Normal}(\mathbf{0},\mathbf{M})$\;
    Simultaneamente amostre $\boldsymbol{\theta}$ e $\boldsymbol{\phi}$ com $L$ \textit{leapfrog steps} cada um reduzido por um fator $\epsilon$.\;
    \Para{$1, 2, \dots, L$}{
     Use o gradiente do $\log$ da posterior de $\boldsymbol{\theta}$ para produzir um \textit{half-step} de $\boldsymbol{\phi}$:
     $\boldsymbol{\phi} \leftarrow \boldsymbol{\phi} + \frac{1}{2} \epsilon \frac{d \log P(\boldsymbol{\theta} \mid y)}{d \theta}$\;
     Use $\boldsymbol{\phi}$ para atualizar $\boldsymbol{\theta}$:
     $\boldsymbol{\theta} \leftarrow \boldsymbol{\theta} + \epsilon \mathbf{M}^{-1} \boldsymbol{\phi}$\;
     Novamente use o gradiente de $\boldsymbol{\theta}$ para produzir um \textit{half-step} de $\boldsymbol{\phi}$:
     $\boldsymbol{\phi} \leftarrow \boldsymbol{\phi} + \frac{1}{2} \epsilon \frac{d \log P(\boldsymbol{\theta} \mid y)}{d \theta}$\;
    }
    Como regra de aceitação/rejeição calcule:
    $r = \frac{P \left(\boldsymbol{\theta}^{(*)} \mid y \right) P \left(\boldsymbol{\phi}^{(*)} \right)}{P \left(\boldsymbol{\theta}^{(t-1)} \mid y \right) P \left(\boldsymbol{\phi}^{(t-1)} \right)}$\;
    Designe:
      $
        \boldsymbol{\theta}^{(t)} =
          \begin{cases}
          \boldsymbol{\theta}^{(*)} & \text{com probabilidade $\min(r,1)$}\\
          \boldsymbol{\theta}^{(t-1)} & \text{caso contrário}
        \end{cases}
      $\;
    \caption{Hamiltonian Monte Carlo (HMC)}
    \end{algorithm}
\end{frame}

\subsubsection{No-U-Turn-Sampler (NUTS)}

\subsubsubsection{As Limitações de HMC e NUTS}

\subsection{Convergência de Correntes Markov}

\subsubsection{Métricas de Convergência}

\subsubsection{Visualizações de Convergência}

\subsubsection{O que fazer se Correntes Markov não Convergirem}
